\section{Exact endpoint and convergent-coordinate tables}
\label{tables:sec}

\subsection{Endpoint corrections through weight eight}
\label{tables:boundaries}

For the unique primitives $H_m=\HH(\A_m)$ in
\eqref{gamma:decomp}, the endpoint functional is
\begin{align*}
 \Delta H_2={}&0,\\
 \Delta H_3={}&-3\zeta(2),\\
 \Delta H_4={}&12\gamma\zeta(2)-10\zeta(3),\\
 \Delta H_5={}&-30\gamma^2\zeta(2)+50\gamma\zeta(3)
                                      +\tfrac{35}{2}\zeta(4),\\
 \Delta H_6={}&60\gamma^3\zeta(2)-150\gamma^2\zeta(3)
 -105\gamma\zeta(4)+210\zeta(2)\zeta(3)-242\zeta(5),\\
 \Delta H_7={}&-105\gamma^4\zeta(2)+350\gamma^3\zeta(3)
                     +\tfrac{735}{2}\gamma^2\zeta(4)\\
 &-1470\gamma\zeta(2)\zeta(3)+1694\gamma\zeta(5)
                +469\zeta(3)^2-\tfrac{1085}{2}\zeta(6),\\
 \Delta H_8={}&168\gamma^5\zeta(2)-700\gamma^4\zeta(3)
                          -980\gamma^3\zeta(4)\\
 &+5880\gamma^2\zeta(2)\zeta(3)-6776\gamma^2\zeta(5)
             -3752\gamma\zeta(3)^2+4340\gamma\zeta(6)\\
 &+7840\zeta(2)\zeta(5)-1862\zeta(3)\zeta(4)-11824\zeta(7).
\end{align*}
They use only the elementary simplifications
$\zeta(2)^2=5\zeta(4)/2$,
$\zeta(2)^3=35\zeta(6)/8$, and
$\zeta(2)\zeta(4)=7\zeta(6)/4$.
The data file \texttt{inverse\_gamma\_normal\_forms.json} retains
the original independent-symbol expressions as well as each exact
primitive and normal form.  Repeating the finite Laurent coefficient
extraction proves every entry of this table.

\subsection{A fully convergent 17-term representation of \texorpdfstring{$\mathcal D(s)$}{D(s)}}
\label{tables:omega-convergent}

The first-index-one table in Section~\ref{sixth:sec} is convenient for
formal harmonic-word generation, but individual entries there may
have spectral poles.  The following alternative avoids them entirely.
The function $\mathcal D(s)$ is the sum of each function in the left
column multiplied by its coefficient in the right column:
\begin{center}
\begin{tabular}{@{}ll@{}}
\toprule
Individually convergent MZV germ & Coefficient \\
\midrule
$\zeta(2+s)$ & $2\zeta(3)-4\gamma\zeta(2)$\\
$\zeta(3+s)$ & $4\gamma^2-8\zeta(2)$\\
$\zeta(4+s)$ & $8\gamma$\\
$\zeta(5+s)$ & $6$\\
$\zeta(2+s,1)$ & $4\zeta(2)$\\
$\zeta(2+s,2)$ & $-4\gamma$\\
$\zeta(2+s,3)$ & $2$\\
$\zeta(3+s,1)$ & $-8\gamma$\\
$\zeta(4+s,1)$ & $-8$\\
$\zeta(2+s,1,2)$ & $4$\\
$\zeta(2+s,2,1)$ & $4$\\
$\zeta(3+s,1,1)$ & $8$\\
$\zeta(2,1+s)$ & $-6\zeta(2)$\\
$\zeta(3,1+s)$ & $8\gamma$\\
$\zeta(4,1+s)$ & $10$\\
$\zeta(2,1+s,2)$ & $-2$\\
$\zeta(3,1+s,1)$ & $-8$\\
\bottomrule
\end{tabular}
\end{center}
Every entry converges absolutely and defines a holomorphic function
on $\re s>-1$.  Therefore $\Omega_5$ is also the sum of the
\emph{individual ordinary derivatives at zero} of these 17 germs,
with the same coefficients.  Derivatives in the second coordinate
must remain second-coordinate derivatives.

For an exact derivation, use the elementary identity, valid for
$k\ge2$ and a positive-index word $\boldsymbol a$,
\begin{equation}\label{tables:upper-tail}
 \begin{split}
 &\sum_{n\ge1}\frac{(H_{n-1}^{(k)}-\zeta(k))
                             H_{n-1}(\boldsymbol a)}{n^{1+s}}\\
 &\hspace{18mm}=-\zeta(k,1+s,\boldsymbol a)
                -\zeta(k+1+s,\boldsymbol a).
 \end{split}
\end{equation}
It follows by writing the negative tail as
$-\sum_{j\ge n}j^{-k}$ and partitioning $j>n$ from $j=n$.
All terms are absolutely summable in the stated half-plane; any
additional inner unit indices contribute at most logarithms.
The coefficient of $1/n$ in \eqref{sixth:delta-r} factors exactly as
\[
 8(H_{n-1}-\gamma)(H_{n-1}^{(3)}-\zeta(3))
 +2(H_{n-1}^{(2)}-\zeta(2))(H_{n-1}^{(2)}+3\zeta(2))
 -10(H_{n-1}^{(4)}-\zeta(4)).
\]
Apply \eqref{tables:upper-tail} to these three terms and retain
the first-index-at-least-two rows of the original table.
Two $\zeta(3+s,2)$ terms cancel; collecting the single zeta terms
gives precisely the 17 rows above.  This establishes the alternative
without analytic regularization of any individual summand.

The same residue calculation gives an exact additional continuation
statement.  Its first two asymptotic coefficients are
$r_n-r_{n-1}=8\zeta(2)(\log n-1)/n^2+
O(\log^2(n+2)/n^3)$, obtained by substituting the ordinary harmonic
expansions in \eqref{sixth:delta-r}.  Consequently
\[
 \mathcal D(s)+8\zeta(2)\{\zeta'(s+2)+\zeta(s+2)\}
\]
is holomorphic for $\re s>-2$.  In particular the exact principal
part of $\mathcal D$ at $s=-1$ is
\[
 \frac{8\zeta(2)}{(s+1)^2}-\frac{8\zeta(2)}{s+1}.
\]
The asymptotic estimate here is used only to prove that spectral
identity and its pole coefficients.

\subsection{A convergent logarithmic antiderivative for the mixed term}

For $a>0$ and $\re z>0$, the pole table in Section~\ref{sixth:sec}
gives the anchored primitive
\begin{equation}\label{tables:primitive6}
 \begin{split}
 \mathcal P_6(z;a)=\sum_{n\ge0}\biggl\{
 &r_n\bigl(\Log(n+z)-\log(n+a)\bigr)\\
 &+\sum_{j=2}^6\frac{p_j(n)}{1-j}
       \bigl((n+z)^{1-j}-(n+a)^{1-j}\bigr)\biggr\}.
 \end{split}
\end{equation}
It converges normally on compact subsets of the right half-plane,
vanishes at $z=a$, and satisfies
$\partial_z\mathcal P_6(z;a)=\psi(z)^2\psi'(z)^2$.
Thus the mixed sixth coordinate has a fully specified antiderivative
as well as its endpoint finite-part value.  The general primitive
theorem supplies the same construction for all the derivative terms
in the seventh, eighth, and higher moment reductions.
